\section{Related Work}\label{subsec:intro-related}

\subsection{Task Planning for Robotic Fleets}
Task planning in robotic fleets addresses the decomposition of missions into
subtasks and their allocation across robots or teams. Surveys on multi-robot
task allocation (MRTA)~\cite{gini2017multi, torreno2017cooperative}
have established taxonomies of problem structures, including vehicle routing,
scheduling, and coalition formation~\cite{dai2024dynamic}.
Centralized optimization methods such as MILP and heuristic
search~\cite{ji2021multi} generate globally near-optimal
solutions and allow precise encoding of resource and time constraints, but the
combinatorial complexity makes them unsuitable for large fleets or missions with
dynamic updates. Decentralized approaches improve scalability and resilience by
distributing decision-making, such as market-based auctions
\cite{zhao2021market} that trade computational effort for speed, and distributed
optimization frameworks~\cite{fioretto2018distributed} that enable concurrent
negotiation among agents. Approximate methods, including evolutionary
algorithms~\cite{liu2019multi} and genetic search~\cite{martin2021mrta}, reduce
computational overhead and can adapt to heterogeneous capabilities, but they
often sacrifice guarantees of feasibility or temporal optimality.
\improve{Recent efforts emphasize hierarchical decomposition and reactive
coordination to improve scalability under dynamic task streams,
including decomposition-based hierarchical planning~\cite{luo2024decomposition,leahy2021scalable}
and real-time reactive allocation under temporal logic~\cite{chen2025real}.}
Despite these advances, most approaches assume static and fully known tasks at the planning
stage. Efficient solutions that adapt to the continual online task generation,
uncertainties in subtask number and distribution,
and online operator interventions remain open.


\subsection{Complex Tasks as Temporal Logic Formulas}
Temporal logic has emerged as a principled way to specify robotic missions that
combine safety, temporal order, and collaboration requirements. Probabilistic
temporal logics (PCTL)~\cite{lahijanian2015temporal} enable reasoning about
uncertainty in outcomes, while LTL provides a rich language for describing
sequential objectives such as repeated patrolling, collaborative capture, or
persistent surveillance~\cite{kantaros2020stylus,  leahy2021scalable, schillinger2018simultaneous,
luo2021temporal}. Counting LTL (cLTL)~\cite{sahin2019multirobot}
extends this by capturing quantitative requirements, for instance requiring a
minimum number of robots at specific subtasks.
Capability LTL~\cite{li2025task,cardona2024planning} address the time-dependent capabilities of robots
during execution.
Hierarchical LTL is proposed in~\cite{luo2024decomposition} \improve{to
further introduce structured analyses and synthesis to improve efficiency}.
Centralized formulations ensure completeness and optimality through sampling-based search
\cite{kantaros2020stylus}, automaton–system synchronization
\cite{schillinger2018simultaneous}, and MILP encoding of task constraints
\cite{leahy2021scalable, luo2021temporal, sahin2019multirobot, kurtz2021more}, but they face
double-exponential growth with task and fleet size. Decentralized methods
address scalability through local coordination~\cite{lindemann2019coupled},
decision trees~\cite{chen2025real}, partial-order analyses~\cite{liu2024time},
reactive synthesis~\cite{kantaros2022perception}, and hierarchical decoupling~\cite{luo2025simultaneous}.
More recent works emphasize dynamic reactivity by combining logic
specifications with scalable heuristics~\cite{leahy2021scalable,chen2025real, cardona2024planning}.
\improve{As summarized in Table~\ref{table:rw_comparison},
  scalability in complex missions has been addressed through mechanisms
  such as sampling~\cite{kantaros2020stylus}, capacity-based optimization~\cite{leahy2021scalable},
  counting constraints~\cite{sahin2019multirobot}, decision tree~\cite{chen2025real},
  and flow-based construction~\cite{gosrich2025online}.
Nevertheless, most aforementioned work overlook the online interaction with human operators,
and different types of local tasks contained within the temporal mission,
which in itself requires coordination and adaptation.}
Such nested coordination raises key challenges in retaining tractable
computational complexity while enabling coherent adaptation across multiple
spatial, temporal and organizational scales.

\improve{As the most relevant to this work, HULK~\cite{luo2025hulk} facilitates continual task
planning by modeling missions constraints as posets. However, this approach necessitates
reconstruction whenever missions are modified online, limiting its online responsiveness.
Furthermore, it fails to account for the interplay between
high-level coordination and low-level motion feasibility, as well as the
integration of human supervision during execution. The proposed method addresses these
shortcomings by incorporating a human interaction protocol and an online adaptation
algorithm, further integrating dynamic constraints into motion feasibility. Additionally,
it enhances computational efficiency by eliminating the need for poset reconstruction
and employing an automata-guided search instead.}

%======================================
\begin{table}[t!]
  \centering
  \improve{
\caption{\textbf{Comparison with Related Work}}
\label{table:rw_comparison}
\vspace{-0.05in}
\setlength{\tabcolsep}{3pt}
\renewcommand{\arraystretch}{1.5}
\resizebox{\columnwidth}{!}{%
\begin{tabular}{l|ccccccc}
\toprule
\textbf{Method} &
\makecell[c]{\scriptsize Task \\ \scriptsize Specification} &
\makecell[c]{\scriptsize Online \\ \scriptsize Task} &
\makecell[c]{\scriptsize Uncertain/ \\ \scriptsize Unknown \\ \scriptsize Subtasks} &
\makecell[c]{\scriptsize Human \\ \scriptsize Interaction} &
\makecell[c]{\scriptsize Dynamic \\ \scriptsize Constraints} &
\makecell[c]{\scriptsize Local \\ \scriptsize Coordination} \\
\midrule
STyLuS*~\cite{kantaros2020stylus} &
\makecell[c]{\small LTL} &
\makecell[c]{\small $\times$} &
\makecell[c]{\small $\times$} &
\makecell[c]{\small $\times$} &
\makecell[c]{\small $\surd$} &
\makecell[c]{\small $\times$} \\
ScRATCHeS~\cite{leahy2021scalable} &
\makecell[c]{\small CaTL} &
\makecell[c]{\small $\times$} &
\makecell[c]{\small $\times$} &
\makecell[c]{\small $\times$} &
\makecell[c]{\small $\surd$} &
\makecell[c]{\small $\times$} \\
DecTree~\cite{chen2025real} &
\makecell[c]{\small LTL} &
\makecell[c]{\small $\surd$} &
\makecell[c]{\small $\surd$} &
\makecell[c]{\small $\times$} &
\makecell[c]{\small $\times$} &
\makecell[c]{\small $\times$} \\
cLTL$+$~\cite{sahin2019multirobot} &
\makecell[c]{\small cLTL} &
\makecell[c]{\small $\times$} &
\makecell[c]{\small $\times$} &
\makecell[c]{\small $\times$} &
\makecell[c]{\small $\surd$} &
\makecell[c]{\small $\times$} \\
HieraTL~\cite{luo2025simultaneous} &
\makecell[c]{\small hLTL} &
\makecell[c]{\small $\surd$} &
\makecell[c]{\small $\surd$} &
\makecell[c]{\small $\times$} &
\makecell[c]{\small $\times$} &
\makecell[c]{\small $\times$} \\
Flow-based~\cite{gosrich2025online} &
\makecell[c]{\small Ordered} &
\makecell[c]{\small $\times$} &
\makecell[c]{\small $\surd$} &
\makecell[c]{\small $\times$} &
\makecell[c]{\small $\times$} &
\makecell[c]{\small $\times$} \\
HULK~\cite{luo2025hulk} &
\makecell[c]{\small LTL} &
\makecell[c]{\small $\surd$} &
\makecell[c]{\small $\surd$} &
\makecell[c]{\small $\times$} &
\makecell[c]{\small $\surd$} &
\makecell[c]{\small $\surd$} \\
\midrule
\textbf{HECTOR (Ours)} &
\makecell[c]{\small LTL} &
\makecell[c]{\small \textbf{$\surd$}} &
\makecell[c]{\small \textbf{$\surd$}} &
\makecell[c]{\small \textbf{$\surd$}} &
\makecell[c]{\small \textbf{$\surd$}} &
\makecell[c]{\small \textbf{$\surd$}} \\
\bottomrule
\end{tabular}%
}}
\vspace{-3mm}
\end{table}
%======================================


\subsection{Human-fleet Interaction}
\improve{Beyond automated task planning, effective human-swarm collaboration under
uncertainty remains a central bottleneck, especially in deciding when and how
operators should validate and intervene online~\cite{ferreira2021survey,dahiya2023survey}.
In practice, no universal metric reliably predicts
whether an autonomous plan will remain operationally
sound once execution begins, so human expertise is essential for validating
reasoning quality and mission viability.
Human–swarm interaction emphasizes interface design, situational awareness, and
scalable mechanisms for directing collectives~\cite{kolling2016human}.
Mixed-initiative approaches~\cite{gombolay2017computational} dynamically shift
decision authority between humans and autonomy, supporting conflict resolution
and efficient replanning in uncertain missions.
Similar approaches are proposed in~\cite{tumova2014maximally,ulusoy2013optimality}
to compute least violating plans.
Scaled-autonomy frameworks~\cite{swamy2020scaledautonomy} extend this by allowing
robots to modulate the
level of assistance provided to the operator based on task load.
Interactive fleet learning~\cite{hoque2022fleetdagger} integrates demonstration,
online supervision, and data aggregation to refine fleet policies over time.
However, many existing paradigms provide limited support for online
validation and intervention, often restricting operators to static displays or
post-hoc plans. In dynamic settings, this can force a fallback to teleoperation, which scales
poorly with the fleet size due to increasing workload and latency. Consequently,
interaction protocols that enable sparse and high-level oversight while preserving
scalable autonomy remain insufficiently addressed.}
